% Figure and table captions. They describe what is shown, without interpretation.
\newcommand{\capFigScaling}{Multi-query \texttt{pmap} on the matched grid, with one query of the repeated-alphabet sequence per chip, the timed region \texttt{apply}, and all eight chips visible. Each marker is one run. (a) Steady-state throughput against the number of chips, on log axes. Dashed lines show one power-law fit to all 16 points, $4528 \times \mathrm{chips}^{0.9631} \times \mathrm{length}^{-1.5724}$ queries/s, drawn at each length. (b) Parallel efficiency, the speedup over one chip divided by the number of chips. The dotted line marks 100\%. At 1000 residues, vertical bars give the range over the rounding of the records. Table~\ref{tab:grid} gives the values.} % src: results/sweep/scaling_law_data.json -> len(grid), grid[*].length; recomputed: V4.fit_printed_coeffs; upper limit of the parallel efficiency ("100%")
\newcommand{\capTabPlatforms}{Baseline Configuration on Three Platforms, by Session}
\newcommand{\capTabExperiments}{Baseline Configuration and TPU Experiments}
\newcommand{\capTabMultichip}{Ways of Mapping AlphaFold2 onto the Eight-Chip Slice at 118 Residues} % src: recomputed: INPUT.toy_length
\newcommand{\capTabSweeps}{Single-Chip Sweeps}
\newcommand{\capTabGrid}{Throughput, Speedup, and Parallel Efficiency on the Matched Grid}
